\documentclass[12pt]{article}
\usepackage{geometry}
\usepackage{hyperref}
\usepackage{lipsum}
\usepackage{enumitem}
\usepackage{titlesec}
\usepackage{graphicx}

\geometry{a4paper, margin=1in}

% Define custom title format for sections
\titleformat{\section}{\normalfont\Large\bfseries}{\thesection.}{1em}{}
\titleformat{\subsection}{\normalfont\large\bfseries}{\thesubsection.}{1em}{}

\title{\textbf{Real-Time Object Detection with AI Interaction}\\[1ex]
\large Detailed Documentation}
\author{Devansh Verma}
\date{\today}

\begin{document}

\maketitle

\tableofcontents

\section{Introduction}

This document provides a detailed overview of the project that integrates real-time object detection using dual YOLO models with a FastAPI-based FAQ chatbot. The system leverages a webcam feed for object detection and uses the chatbot to deliver relevant information about detected objects. Additionally, it incorporates text-to-speech (TTS) functionality for enhanced interactive user experiences.

\section{Project Overview}

\subsection{Approach \& Methodology}

The project is divided into two primary components:
\begin{enumerate}[label=\arabic*.]
    \item \textbf{Real-Time Object Detection:}
    \begin{itemize}
        \item Utilizes dual YOLO models:
        \begin{itemize}
            \item A custom-trained YOLO model (``best.pt'') tailored for domain-specific objects.
            \item A COCO-pretrained YOLO model (``yolov5n.pt'') for general object detection.
        \end{itemize}
        \item Captures video frames from a webcam, processes them in real-time, and applies concurrent inference to speed up processing.
        \item Uses OpenCV to draw bounding boxes and overlay text information on the video feed.
    \end{itemize}
    
    \item \textbf{FAQ Chatbot API:}
    \begin{itemize}
        \item Built using FastAPI and deployed at the provided URL.
        \item Processes FAQ data stored in a CSV file (``improved\_faq-1.csv'') with two columns: \textit{question} and \textit{answer}.
        \item Uses SentenceTransformers (``all-MiniLM-L6-v2'') to encode FAQ questions into embeddings.
        \item Implements FAISS for efficient similarity search based on the embedding space.
        \item Integrates thresholding to ensure only sufficiently similar questions retrieve an answer.
    \end{itemize}
\end{enumerate}

\subsection{Integration Workflow}

\begin{enumerate}[label=\arabic*.]
    \item The object detection script processes each frame and identifies objects using both YOLO models.
    \item For each detected object (except persons), a query is generated and sent to the FAQ chatbot.
    \item The FAQ chatbot searches its FAISS index for the most similar question and returns the corresponding answer.
    \item Text-to-speech (TTS) is used to audibly deliver the response.
    \item If a person is detected, the system prompts the user to manually enter a query.
\end{enumerate}

\section{Data Preprocessing \& Selection}

\subsection{FAQ Data}

\begin{itemize}
    \item \textbf{Source:} FAQ data is maintained in a CSV file named \texttt{improved\_faq-1.csv}.
    \item \textbf{Structure:} The file contains two columns:
    \begin{itemize}
        \item \textbf{question:} The FAQ question.
        \item \textbf{answer:} The corresponding answer.
    \end{itemize}
    \item \textbf{Preprocessing:} 
    \begin{itemize}
        \item Column headers are standardized.
        \item The questions are encoded into high-dimensional embeddings using the SentenceTransformer model.
    \end{itemize}
\end{itemize}

\subsection{Indexing with FAISS}

\begin{itemize}
    \item The encoded FAQ embeddings are stored in a FAISS index using an L2 distance metric (\texttt{IndexFlatL2}).
    \item The FAISS index allows for rapid similarity searches when processing user queries.
    \item The index is saved to disk as \texttt{faq\_index.faiss} and is reloaded along with the processed FAQ data (\texttt{faq\_data.csv}) during runtime.
\end{itemize}

\section{Model Architecture \& Tuning Process}

\subsection{YOLO Models}

\begin{itemize}
    \item \textbf{Custom YOLO Model (``best.pt''):}
    \begin{itemize}
        \item Trained on a domain-specific dataset (e.g., hotel items such as ``ac'', ``bed'', ``sofa'', etc.).
        \item Fine-tuned with adjustments to learning rate, batch size, and number of epochs based on validation performance.
    \end{itemize}
    \item \textbf{COCO YOLO Model (``yolov5n.pt''):}
    \begin{itemize}
        \item A lightweight variant pre-trained on the COCO dataset.
        \item Provides fast inference speeds to support real-time processing.
    \end{itemize}
\end{itemize}

\subsection{FAQ Chatbot Components}

\begin{itemize}
    \item \textbf{SentenceTransformer:} The model used is \texttt{all-MiniLM-L6-v2}, which provides efficient sentence embeddings.
    \item \textbf{FAISS:} An L2-based flat index is used to search for similar questions based on embedding distances.
    \item \textbf{Thresholding:} A similarity threshold (e.g., \texttt{SIMILARITY\_THRESHOLD = 1.0}) is defined to filter out low-confidence matches.
\end{itemize}

\subsection{Detailed Usage of Dual YOLO Models}

The project employs two distinct YOLO models to maximize the range and accuracy of object detection:

\begin{enumerate}[label=\arabic*.]
    \item \textbf{Custom YOLO Model (``best.pt''):}
    \begin{itemize}
        \item \textbf{Purpose:} This model is specifically trained on domain-relevant objects such as those encountered in hotel environments or Airbnb-style settings. It is optimized to detect items like air conditioners, beds, sofas, and other furniture with high precision.
        \item \textbf{Advantages:}
        \begin{itemize}
            \item \textbf{Domain-Specific Accuracy:} Its training on a tailored dataset allows for a more nuanced detection of objects that are critical for interior layout analysis.
            \item \textbf{Fine-Tuned Detection:} It can identify subtle differences in similar items, ensuring that the responses about room layout and related elements are highly detailed.
        \end{itemize}
        \item \textbf{Usage in the System:} 
        \begin{itemize}
            \item When the custom model detects an object, it provides responses that include detailed contextual information, particularly useful for room layout analysis and Airbnb property descriptions.
        \end{itemize}
    \end{itemize}
    
    \item \textbf{COCO YOLO Model (``yolov5n.pt''):}
    \begin{itemize}
        \item \textbf{Purpose:} This model is pre-trained on the COCO dataset, which contains a broad spectrum of everyday objects. It is integrated to extend the detection capability beyond the specialized domain.
        \item \textbf{Advantages:}
        \begin{itemize}
            \item \textbf{General Object Detection:} Capable of detecting a wide variety of objects such as vehicles, people, animals, and household items.
            \item \textbf{Versatility:} Enables the system to respond to queries related to general objects, thereby increasing the range of interactions.
        \end{itemize}
        \item \textbf{Usage in the System:} 
        \begin{itemize}
            \item By combining the COCO model with the custom model, the system can provide comprehensive responses. If an everyday object is detected (e.g., a car or a pet), the COCO model triggers an appropriate FAQ response.
        \end{itemize}
    \end{itemize}
\end{enumerate}

\textbf{Combined Benefits:}
\begin{itemize}
    \item The dual-model approach ensures high accuracy for both domain-specific objects and general everyday objects.
    \item In cases where both models detect the same object, a merging strategy based on confidence scores selects the optimal detection, thereby minimizing false positives.
    \item This comprehensive detection strategy allows the system to offer richer and more diverse responses, enhancing its utility in varied scenarios.
\end{itemize}

\section{Performance Results \& Next Steps}

\subsection{Performance Results}

\begin{itemize}
    \item \textbf{Object Detection:}
    \begin{itemize}
        \item The combination of the custom and COCO YOLO models yields robust object detection across a range of environments.
        \item Real-time performance is maintained by resizing frames and using concurrent processing.
    \end{itemize}
    \item \textbf{FAQ Chatbot Accuracy:}
    \begin{itemize}
        \item The SentenceTransformer and FAISS integration delivers quick and accurate FAQ responses.
        \item The threshold mechanism helps in reducing false positives in FAQ retrieval.
    \end{itemize}
\end{itemize}

\subsection{Next Steps \& Future Improvements}

\begin{itemize}
    \item \textbf{Model Optimization:} Further fine-tuning of the custom YOLO model with additional domain-specific data.
    \item \textbf{Scalability:} Deploy the FAQ chatbot on a scalable cloud platform to handle higher loads and reduce latency.
    \item \textbf{Enhanced Interaction:} Integrate voice recognition for hands-free operation and implement a more sophisticated dialog management system.
    \item \textbf{User Feedback:} Develop a feedback loop for users to rate responses, enabling continuous improvement of both detection and response quality.
\end{itemize}

\section{Integration Potential}

This project presents significant potential for integration with various systems and applications, including:

\begin{itemize}
    \item \textbf{Smart Home and IoT:} The real-time object detection system can be integrated into smart home ecosystems to monitor room environments, trigger alerts, or adjust settings based on detected objects. For example, detecting the presence of a person could automatically adjust lighting or climate control.
    \item \textbf{Retail and Hospitality:} In environments such as hotels or rental properties, the system can be integrated into property management solutions to provide guests with immediate information about room amenities and local attractions.
    \item \textbf{Security and Surveillance:} The dual-model approach enhances security systems by providing precise detection of both domain-specific items and general objects. Integration with surveillance systems can lead to improved anomaly detection and automated alert systems.
    \item \textbf{Healthcare and Assisted Living:} The system’s ability to detect objects in real-time can be applied to assistive technologies in healthcare, such as monitoring patient activity or ensuring safety in assisted living facilities.
    \item \textbf{Data Analytics and Business Intelligence:} Integration with analytics platforms can allow for the aggregation of detection data to gain insights into user interactions and object frequency, which can then inform operational decisions and targeted improvements.
\end{itemize}

\section{Conclusion}

This project demonstrates a successful integration of real-time object detection with an FAQ chatbot, leveraging the complementary strengths of a custom YOLO model and a COCO-pretrained YOLO model. The dual-model approach not only enhances domain-specific accuracy but also broadens the detection capability to include a wide range of everyday objects. Combined with efficient FAQ retrieval and TTS, the system delivers a rich interactive experience. Contributions and feedback are welcome.

\section*{Acknowledgments}

\begin{itemize}
    \item \textbf{YOLO Models:} Special thanks to Ultralytics for their YOLO implementations.
    \item \textbf{FAISS:} Developed by Facebook AI Research, FAISS enables efficient similarity search.
    \item \textbf{SentenceTransformers:} Provides state-of-the-art sentence embeddings.
    \item \textbf{FastAPI:} For rapid development of the FAQ chatbot API.
    \item \textbf{OpenCV:} Essential for real-time video processing.
\end{itemize}

\end{document}
